\scanpage{71}
\mspage{20}

Dans la théorie de $\mathrm{HOT}_W$, influencé par la réflexion générale sur les dérivateurs et l'axiome Der 3 b, j'ai travaillé sous l'hypothèse $W\supset W_\infty$, mais en utilisant un axiome de fibration $\mathrm L\,5=W(4\,\mathrm b)$ qui, il s'avère, implique fatalement $W=W_\infty$. J'ai l'impression que je n'ai jusqu'à présent (propriétés de $f_!$ pour $\mathrm{HOT}$) pas utilisé $\unsure{V|W}\supset W_\infty$, donc ça peut être utile, comme ``ça ne coûte pas plus cher'' (et oblige d'être plus attentif à ce qu'on utilise), de laisser tomber $W\supset W_\infty$ \deleted{\unsure{et}} \deleted{\unsure{de de}} (impliquant $W=W_\infty$), et de mettre par contre sur $W$ des axiomes aussi forts (concernant \deleted{\unsure{sa}} fibration) que les propriétés de $W_\infty$, ``sans se priver en rien''. Je supposerai donc $W(1)$, $W(2)$, $W(3)$, $W(4\,\mathrm a)$ (la forme la plus forte de l'axiome des filtrations (cf pages 13, 14) -- c'est tout ce qui a été \unsure{utile|utilisé} jusqu'à présent. \deleted{Je} Si on y joint $W(5)$ (axiome de connexité) ça implique $W\subset W_\infty$. Mais pour l'instant je ne \unsure{crois|vois} pas en avoir un besoin. \deleted{Mais}

Mais méditant sur la question, je crois avoir trouvé le
% slanted
\margin{Cf pour un résultat plus fort XV \deleted{\illegible{2}} 9.12 p. \unsure{133bis3}.} \uline{Théorème}? (!) : Soit $W\subset\mathrm{Fl}(\mathrm{Cat})$ satisfaisant $W(1\deleted{\unsure{\mathrm b}},2,3,4\,\mathrm b'')$. Si $W$ ne satisfait pas à $W(5)$, alors $W=\mathrm{Fl}(\mathrm{Cat})$.
% slanted
\margin{\uline{NB} Le résultat \unsure{l.c.|cité} est plus général mais pas vraiment plus fort, car il ne donne pas $W=\mathrm{Fl}(\mathrm{Cat})$} \uline{Dém} : Soit $f:X\longrightarrow Y$ de $\mathrm{Cat}$, tel que $\pi_0(f)$ ne soit pas surjectif. Soit $I$ la catégorie discrète ayant pour \inserted{ens. d'}objets l'ens $\pi_0(Y)$, on a donc un diagramme commutatif \diagram{\begin{tikzcd}[ampersand replacement=\&,column sep=small,row sep=small] X\arrow[rr,"f"]\arrow[dr,"p"']\&\&Y\arrow[dl,"q"]\\ \& I\& \end{tikzcd}} d'autre part \inserted{un foncteur de} une catégorie \deleted{dans} une catégorie discrète est toujours $W_\infty$-parfait, donc $p,q$ sont des $W$-fibrations. Comme

% archive sheet number: 71
